% GENERATED FILE. Edit canonical lessons, then run npm run generate:books.
% canonical-source-hash: fnv1a64:89bd51735977d9aa
% canonical-lessons: GE-C345-zusammenhang, GE-C345-zusammenhaengen, GE-C345-beweis, GE-C345-belegen, GE-C345-nachweisen, GE-R345-review-338-345

\chapter{Zusammenhänge}
\label{ch:ge-zusammenhange}
\hlchaptermodality{345}

\begin{chapteropening}
\textbf{By the end of this chapter:} \emph{I can say that two things are connected, give proof and back up a claim.}

Complete the last lesson of chapter 345: I can say that two things are connected, give proof and back up a claim.
\end{chapteropening}

\section[der Zusammenhang]{der Zusammenhang — the connection, the context}

\label{lesson:GE-C345-zusammenhang}



\begin{quote}
Before the new one: say the German for to affect, then the German for to concern.
\end{quote}

\subsection*{You'll want to know: der Zusammenhang}
\textbf{der Zusammenhang} — \textquotedblleft{}the connection, the context\textquotedblright{}.

\begin{quote}
Es gibt einen Zusammenhang zwischen Schlaf und Gesundheit.
\end{quote}

\begin{quote}
\emph{There is a connection between sleep and health.}
\end{quote}

Plural: \emph{die Zusammenhänge}. \emph{In diesem Zusammenhang} — in this context.

\begin{cousinweb}
\emph{Zusammenhang} is built from \emph{zusammen} \textquotedblleft{}together\textquotedblright{} and the family of \emph{hängen} \textquotedblleft{}to hang\textquotedblright{}, from Old High German \emph{hengen} and Proto-Germanic \emph{hangijaną} (EtymDB): things that hang together.
\end{cousinweb}

\begin{grammarlens}[title={What you've built}]
One more word for telling, explaining and describing at B1.
\end{grammarlens}

\subsection*{Your turn}
\begin{itemize}
\raggedright
  \item \emph{Say it:} \emph{der Zusammenhang}
  \item \emph{Say it:} \emph{der Zusammenhang}, once more
  \item \emph{Say it:} a sentence of your own with \emph{der Zusammenhang}
  \item \emph{Recall:} say the German for to affect, then the German for to concern, then say \emph{der Zusammenhang} again
\end{itemize}

\begin{tcolorbox}[breakable,colback=teal!4,colframe=teal!35!black,title={Before you move on}]
What does der Zusammenhang mean? (\textquotedblleft{}The connection, the context\textquotedblright{}.) Say it once more.
\end{tcolorbox}
\section[zusammenhängen]{zusammenhängen — to be connected}

\label{lesson:GE-C345-zusammenhaengen}



\begin{quote}
Before the new one: say the German for to concern, then the German for the connection.
\end{quote}

\subsection*{You'll want to know: zusammenhängen}
\textbf{zusammenhängen} — \textquotedblleft{}to be connected\textquotedblright{}.

\begin{quote}
Seine schlechte Laune hängt mit der Arbeit zusammen.
\end{quote}

\begin{quote}
\emph{His bad mood is connected with his work.}
\end{quote}

Separable and strong: \emph{das hängt mit … zusammen}; Perfekt \emph{hat zusammengehangen}.

\begin{cousinweb}
\emph{zusammenhängen} is built from \emph{zusammen} \textquotedblleft{}together\textquotedblright{} and \emph{hängen} \textquotedblleft{}to hang\textquotedblright{}, from Old High German \emph{hengen} and Proto-Germanic \emph{hangijaną} (EtymDB).
\end{cousinweb}

\begin{grammarlens}[title={What you've built}]
One more word for telling, explaining and describing at B1.
\end{grammarlens}

\subsection*{Your turn}
\begin{itemize}
\raggedright
  \item \emph{Say it:} \emph{zusammenhängen}
  \item \emph{Say it:} \emph{zusammenhängen}, once more
  \item \emph{Say it:} a sentence of your own with \emph{zusammenhängen}
  \item \emph{Recall:} say the German for to concern, then the German for the connection, then say \emph{zusammenhängen} again
\end{itemize}

\begin{tcolorbox}[breakable,colback=teal!4,colframe=teal!35!black,title={Before you move on}]
What does zusammenhängen mean? (\textquotedblleft{}To be connected\textquotedblright{}.) Say it once more.
\end{tcolorbox}
\section[der Beweis]{der Beweis — the proof, the evidence}

\label{lesson:GE-C345-beweis}



\begin{quote}
Before the new one: say the German for the connection, then the German for to be connected.
\end{quote}

\subsection*{You'll want to know: der Beweis}
\textbf{der Beweis} — \textquotedblleft{}the proof, the evidence\textquotedblright{}.

\begin{quote}
Die Polizei hat keinen Beweis gegen den Mann.
\end{quote}

\begin{quote}
\emph{The police have no proof against the man.}
\end{quote}

Plural: \emph{die Beweise}. \emph{Der Beweis für} + accusative. The verb is \emph{beweisen} — to prove.

\begin{cousinweb}
\emph{Beweis} is the noun of \emph{beweisen} \textquotedblleft{}to prove\textquotedblright{}, built with \emph{be-} on \emph{weisen} \textquotedblleft{}to show, to direct\textquotedblright{}: a proof shows the thing. EtymDB has no history for the noun, so this is the plain make-up.
\end{cousinweb}

\begin{grammarlens}[title={What you've built}]
One more word for telling, explaining and describing at B1.
\end{grammarlens}

\subsection*{Your turn}
\begin{itemize}
\raggedright
  \item \emph{Say it:} \emph{der Beweis}
  \item \emph{Say it:} \emph{der Beweis}, once more
  \item \emph{Say it:} a sentence of your own with \emph{der Beweis}
  \item \emph{Recall:} say the German for the connection, then the German for to be connected, then say \emph{der Beweis} again
\end{itemize}

\begin{tcolorbox}[breakable,colback=teal!4,colframe=teal!35!black,title={Before you move on}]
What does der Beweis mean? (\textquotedblleft{}The proof, the evidence\textquotedblright{}.) Say it once more.
\end{tcolorbox}
\section[belegen]{belegen — to back up (with evidence); to cover}

\label{lesson:GE-C345-belegen}



\begin{quote}
Before the new one: say the German for to be connected, then the German for the proof.
\end{quote}

\subsection*{You'll want to know: belegen}
\textbf{belegen} — \textquotedblleft{}to back up (with evidence); to cover\textquotedblright{}.

\begin{quote}
Die Zahlen belegen, dass mehr Menschen Fahrrad fahren.
\end{quote}

\begin{quote}
\emph{The figures back up the fact that more people are cycling.}
\end{quote}

Inseparable: \emph{es belegt, es hat belegt}. Also to cover: \emph{ein belegtes Brötchen} — a filled roll; \emph{der Platz ist belegt} — the seat is taken.

\begin{cousinweb}
\emph{belegen} is inherited from Old High German \emph{bileggen} (EtymDB), built with \emph{be-} on \emph{legen} \textquotedblleft{}to lay\textquotedblright{}: you lay proof down. EtymDB lists English \emph{belay} as a cousin.
\end{cousinweb}

\begin{grammarlens}[title={What you've built}]
One more word for telling, explaining and describing at B1.
\end{grammarlens}

\subsection*{Your turn}
\begin{itemize}
\raggedright
  \item \emph{Say it:} \emph{belegen}
  \item \emph{Say it:} \emph{belegen}, once more
  \item \emph{Say it:} a sentence of your own with \emph{belegen}
  \item \emph{Recall:} say the German for to be connected, then the German for the proof, then say \emph{belegen} again
\end{itemize}

\begin{tcolorbox}[breakable,colback=teal!4,colframe=teal!35!black,title={Before you move on}]
What does belegen mean? (\textquotedblleft{}To back up (with evidence); to cover\textquotedblright{}.) Say it once more.
\end{tcolorbox}
\section[nachweisen]{nachweisen — to prove, to demonstrate}

\label{lesson:GE-C345-nachweisen}



\begin{quote}
Before the new one: say the German for the proof, then the German for to back up.
\end{quote}

\subsection*{You'll want to know: nachweisen}
\textbf{nachweisen} — \textquotedblleft{}to prove, to demonstrate\textquotedblright{}.

\begin{quote}
Der Mann konnte nachweisen, dass er zu Hause war.
\end{quote}

\begin{quote}
\emph{The man was able to prove that he was at home.}
\end{quote}

Separable and strong: \emph{er weist nach}, \emph{er hat nachgewiesen}.

\begin{cousinweb}
\emph{nachweisen} is built from \emph{nach} and \emph{weisen} \textquotedblleft{}to show, to direct\textquotedblright{} (EtymDB): you show the way back to the facts.
\end{cousinweb}

\begin{grammarlens}[title={What you've built}]
One more word for telling, explaining and describing at B1.
\end{grammarlens}

\subsection*{Your turn}
\begin{itemize}
\raggedright
  \item \emph{Say it:} \emph{nachweisen}
  \item \emph{Say it:} \emph{nachweisen}, once more
  \item \emph{Say it:} a sentence of your own with \emph{nachweisen}
  \item \emph{Recall:} say the German for the proof, then the German for to back up, then say \emph{nachweisen} again
\end{itemize}

\begin{tcolorbox}[breakable,colback=teal!4,colframe=teal!35!black,title={Before you move on}]
What does nachweisen mean? (\textquotedblleft{}To prove, to demonstrate\textquotedblright{}.) Say it once more.
\end{tcolorbox}
\section[(dialogue)]{Words from chapters 338–345 — a review}

\label{lesson:GE-R345-review-338-345}



\begin{quote}
Say the German for the guess and the assumption.
\end{quote}

\subsection*{Your turn}
Say these aloud:

\begin{itemize}
\raggedright
  \item the guess, the assumption, apparently, obviously, to suspect
  \item the remark, to indicate, the hint, to bring out, to set out
  \item to judge, to assess, the assessment, to criticize, the criticism
  \item although, admittedly, however, nevertheless, in fact
  \item the effect, to work, to cause, to trigger, the trigger
  \item the consequence, the result, consequently, so that, therefore
  \item to influence, the influence, the impact, to affect, to concern
  \item the connection, to be connected, the proof, to back up, to prove
\end{itemize}

\begin{tcolorbox}[breakable,colback=teal!4,colframe=teal!35!black,title={Before you move on}]
The guess? (\textbf{die Vermutung}.) The assumption? (\textbf{die Annahme}.) Apparently? (\textbf{anscheinend}.) Obviously? (\textbf{offensichtlich}.) To suspect? (\textbf{ahnen}.) The remark? (\textbf{die Äußerung}.) To indicate? (\textbf{hinweisen}.) The hint? (\textbf{der Hinweis}.) To bring out? (\textbf{hervorheben}.) To set out? (\textbf{darlegen}.) To judge? (\textbf{beurteilen}.) To assess? (\textbf{einschätzen}.) The assessment? (\textbf{die Einschätzung}.) To criticize? (\textbf{kritisieren}.) The criticism? (\textbf{die Kritik}.) Although? (\textbf{obwohl}.) Admittedly? (\textbf{zwar}.) However? (\textbf{allerdings}.) Nevertheless? (\textbf{dennoch}.) In fact? (\textbf{tatsächlich}.) The effect? (\textbf{die Wirkung}.) To work? (\textbf{wirken}.) To cause? (\textbf{verursachen}.) To trigger? (\textbf{auslösen}.) The trigger? (\textbf{der Auslöser}.) The consequence? (\textbf{die Konsequenz}.) The result? (\textbf{das Resultat}.) Consequently? (\textbf{folglich}.) So that? (\textbf{sodass}.) Therefore? (\textbf{daher}.) To influence? (\textbf{beeinflussen}.) The influence? (\textbf{der Einfluss}.) The impact? (\textbf{die Auswirkung}.) To affect? (\textbf{auswirken}.) To concern? (\textbf{betreffen}.) The connection? (\textbf{der Zusammenhang}.) To be connected? (\textbf{zusammenhängen}.) The proof? (\textbf{der Beweis}.) To back up? (\textbf{belegen}.) To prove? (\textbf{nachweisen}.)
\end{tcolorbox}
